\documentclass[12pt,a4paper]{article}

% Paket
\usepackage[utf8]{inputenc}
\usepackage[T1]{fontenc}
\usepackage[swedish]{babel}
\usepackage[autostyle=true]{csquotes}
\usepackage[margin=2.5cm]{geometry}
\usepackage{graphicx}
\usepackage{xcolor}
\usepackage{hyperref}
\usepackage{amsmath}
\usepackage{booktabs}
\usepackage{enumitem}
\usepackage{listings}

% Undvik konflikt med babels "-shorthand i löpande text
\AtBeginDocument{\shorthandoff{"}}

% Definiera JSON för listings
\lstdefinelanguage{json}{
  morestring=[b]",
  morestring=[d]'
}

% Färger
\definecolor{brandgreen}{RGB}{0,112,74}
\definecolor{lightgray}{RGB}{245,245,245}

% Hyperref
\hypersetup{
  colorlinks=true,
  linkcolor=brandgreen,
  urlcolor=brandgreen,
  citecolor=brandgreen
}

% Titel
\title{\textbf{Metod för riskbedömning av kund} \\ \large Version 3.0 – Celestial Risk Engine}
\author{Celestial AB \\ Penningtvättsansvarig}
\date{\today}

\begin{document}

\maketitle
\tableofcontents
\newpage

% ========================================================================
\section{Introduktion och metodval}
% ========================================================================

\subsection{Syfte}

Detta dokument beskriver Celestial AB:s metod för riskbedömning av kunder enligt Penningtvättslagen (PTL 2017:630). Metoden används för att klassificera nya kunder i riskklasser och avgöra om Enhanced Due Diligence (EDD) krävs.

\textbf{Version 3.0 ersätter tidigare Hot×Sårbarhet-modell} och inför en enhetlig 0–100 poäng-baserad Celestial Risk Engine.

\subsection{Regelverkskrav}

\begin{itemize}[noitemsep]
  \item \textbf{PTL 3 kap 2 §:} Riskbaserade kundkännedomsåtgärder
  \item \textbf{PTL 3 kap 10 §:} Fördjupad kundkännedom vid förhöjd risk
  \item \textbf{FFFS 2009:1:} Finansinspektionens föreskrifter om åtgärder mot penningtvätt och finansiering av terrorism
  \item \textbf{EU-direktiv 2015/849 (4:e penningtvättsdirektivet):} Riskbaserat arbetssätt
\end{itemize}

\subsection{Metodöversikt}

\textbf{Celestial Risk Engine} kombinerar tre datakällor:

\begin{enumerate}[noitemsep]
  \item \textbf{Offentliga register} (30\% vikt): Bolagsverket, SPAR, Roaring (PEP, sanktioner, näringsförbud)
  \item \textbf{Korsvalidering} (30\% vikt): Avvikelser mellan registerdata (VH-mismatch, adressmismatch, F-skatt/moms)
  \item \textbf{Bokföringsdata} (40\% vikt): SIE-fil analys, bankkonto, nyckeltal
\end{enumerate}

\textbf{Steg 1 (Statisk KYC)} beräknar en \texttt{staticKycScore} (0–100) baserat på punkt 1–2 ovan (bokföring = 0 initialt).

\textbf{Senare steg} lägger till bokföringskomponenten (40\%) när SIE-fil och bankkonto analyseras.

% ========================================================================
\section{SNI-kod och Hotfaktor}
% ========================================================================

\subsection{Hotfaktor (1–4) från branschkod}

Varje företags SNI-kod (Svensk Näringsgrensindelning) mappas till en \textbf{Hotfaktor} enligt Finansinspektionens riskklassificering av branscher:

\begin{table}[h]
\centering
\begin{tabular}{cl}
\toprule
\textbf{Hot} & \textbf{Branschtyp} \\
\midrule
1 & Lågrisk: Konsulttjänster, IT, utbildning, design \\
2 & Normal risk: Handel (ej kontanter), tillverkning, bygg \\
3 & Förhöjd risk: Restaurang, taxi, frisör, bilhandel (kontantintensiv) \\
4 & Hög risk: Spelverksamhet, valutaväxling, ädelmetaller, begagnathandel \\
\bottomrule
\end{tabular}
\caption{SNI-kod till Hotfaktor-mappning}
\end{table}

\textbf{Källa:} Bolagsverket API returnerar SNI-kod automatiskt vid validering av organisationsnummer.

\subsection{Hotfaktorns påverkan på riskscore}

Hotfaktorn används för att lägga till en \textbf{hotbonus} direkt till den statiska KYC-scoren:

\begin{equation}
\text{hotBonus} = (\text{sniHotFactor} - 1) \times 10
\end{equation}

\begin{equation}
\text{adjustedScore} = \text{publicData} + \text{crossValidation} + \text{bookkeeping} + \text{hotBonus}
\end{equation}

\textbf{Exempel:}
\begin{itemize}[noitemsep]
  \item Hot 1 (lågrisk): hotBonus = $(1-1) \times 10 = 0$ poäng
  \item Hot 2 (normal): hotBonus = $(2-1) \times 10 = 10$ poäng
  \item Hot 3 (förhöjd): hotBonus = $(3-1) \times 10 = 20$ poäng
  \item Hot 4 (hög): hotBonus = $(4-1) \times 10 = 30$ poäng
\end{itemize}

\textbf{Motivering:} Kontantintensiva branscher (Hot 3–4) kräver extra granskning enligt PTL 3 kap 10 § 2 p. Den direkta additionen gör beräkningen transparent och lätt att följa i granskningssyfte.

% ========================================================================
\section{Celestial Risk Engine – Beräkningsmodell}
% ========================================================================

\subsection{Komponent 1: Offentliga register (max 30 poäng)}

\textbf{Datakällor:}
\begin{itemize}[noitemsep]
  \item Bolagsverket: Företagsålder, organisationsform, status
  \item SPAR: Folkbokföring VD/firmatecknare
  \item Roaring: PEP-status, sanktionslistor, näringsförbud, verkliga huvudmän
\end{itemize}

\textbf{Poängberäkning:}

\begin{table}[h]
\centering
\begin{tabular}{lcc}
\toprule
\textbf{Faktor} & \textbf{Villkor} & \textbf{Poäng} \\
\midrule
Bolagsålder & $< 1$ år & +8 \\
            & 1–3 år & +4 \\
            & $> 3$ år & 0 \\
PEP-status  & Någon huvudman är PEP & +10 \\
Sanktioner  & Träff på sanktionslista & +30 (→ critical) \\
Näringsförbud & Aktiv näringsförbudad person & +30 (→ critical) \\
Utländsk VD & VD ej folkbokförd i Sverige & +5 \\
\bottomrule
\end{tabular}
\caption{Offentliga register – poängtabell}
\end{table}

\textbf{Max-tak:} 30 poäng (utom vid sanktioner/näringsförbud som ger automatisk \texttt{critical}-klassning).

\subsection{Komponent 2: Korsvalidering (max 30 poäng)}

\textbf{Avvikelser mellan registerdata:}

\begin{table}[h]
\centering
\begin{tabular}{lcc}
\toprule
\textbf{Avvikelse} & \textbf{Beskrivning} & \textbf{Poäng} \\
\midrule
VH-mismatch & Användare påstår sig vara VH men Roaring visar ej & +15 \\
Adressmismatch & SPAR-adress $\neq$ Bolagsverket-adress & +8 \\
Saknad F-skatt & Fakturerar men ej F-skattesedel & +10 \\
Saknad momsreg. & Omsättning $> 30\,000$ kr men ej momsregistrerad & +12 \\
Dubblerad VD & Samma person VD i $>5$ bolag & +6 \\
\bottomrule
\end{tabular}
\caption{Korsvalidering – poängtabell}
\end{table}

\textbf{Max-tak:} 30 poäng.

\textbf{Motivering:} Avvikelser mellan register indikerar potentiella strukturer för skatteundandragande eller penningtvätt (PTL 3 kap 10 § 1 p).

\subsection{Komponent 3: Bokföringsdata (max 40 poäng)}

\textbf{Analyseras i senare steg} när SIE-fil och bankkonto laddas upp.

\textbf{Indikatorer:}
\begin{itemize}[noitemsep]
  \item Kassaandel $> 40\%$ av omsättning (+12 poäng)
  \item Negativ soliditet $< -20\%$ (+8 poäng)
  \item Extraordinära transaktioner $> 500\,000$ kr utan verifikat (+15 poäng)
  \item Utländska transaktioner till högrisländer (FATF grå-/svartlista) (+10 poäng)
  \item Avvikelse bokföring vs. deklaration $> 15\%$ (+8 poäng)
\end{itemize}

\textbf{Max-tak:} 40 poäng.

\textbf{I Steg 1 sätts denna komponent till 0}, eftersom bokföring ännu ej analyserats.

% ========================================================================
\section{Riskklassificering – Fyrgradig skala}
% ========================================================================

\subsection{Finansinspektionens standard}

Celestial AB använder Finansinspektionens fyrgradig riskskala (FFFS 2009:1):

\begin{enumerate}[noitemsep]
  \item \textbf{Låg risk} (Low): Standardiserad kundkännedom
  \item \textbf{Förhöjd risk} (Medium): Utökad kundkännedom
  \item \textbf{Hög risk} (High): Enhanced Due Diligence (EDD) enligt PTL 3 kap 10 §
  \item \textbf{Oacceptabel risk} (Critical): Avslag av uppdrag, anmälningsplikt Finanspolisen
\end{enumerate}

\subsection{Tröskelvärden}

\textbf{Beräkning:}
\begin{equation}
\text{adjustedScore} = \text{publicData} + \text{crossValidation} + \text{bookkeeping} + \text{hotBonus}
\end{equation}

där $\text{hotBonus} = (\text{sniHotFactor} - 1) \times 10$

\begin{table}[h]
\centering
\begin{tabular}{lcc}
\toprule
\textbf{Riskklass} & \textbf{adjustedScore} & \textbf{Åtgärd} \\
\midrule
Low & 0–40 poäng & Standardiserad KYC \\
Medium & 41–70 poäng & Utökad KYC, extra dokumentation \\
High & 71–90 poäng & EDD, godkännande av PTL-ansvarig \\
Critical & $> 90$ poäng eller hårda triggers & Avslag, eventuell anmälan \\
\bottomrule
\end{tabular}
\caption{Riskklassificering – tröskelvärden}
\end{table}

\subsection{Hårda triggers (automatisk Critical)}

Oavsett score klassas kunden som \texttt{critical} vid:
\begin{itemize}[noitemsep]
  \item Träff på sanktionslista (Roaring Sanctions API)
  \item Aktivt näringsförbud (Roaring Business Prohibition API)
  \item Flera allvarliga avvikelser: VH-mismatch + saknad F-skatt + PEP
\end{itemize}

\subsection{Eskaleringstriggrar (höjer riskklass ett steg)}

Om \texttt{adjustedScore} ligger på gränsen (t.ex. 38–42 eller 68–72) eskaleras riskklass vid:
\begin{itemize}[noitemsep]
  \item PEP-status + score $> 35$ → Low blir Medium
  \item VH-mismatch + score $> 60$ → Medium blir High
  \item Utländska partners (högrisländer) + score $> 55$ → Medium blir High
\end{itemize}

% ========================================================================
\section{Implementering i Steg 1 API}
% ========================================================================

\subsection{Response-struktur}

\textbf{Endpoint:} \texttt{POST /api/onboarding/\{onboardingId\}/riskfragor/steg1}

\textbf{Response (förenklad):}
\begin{lstlisting}[language=json]
{
  "success": true,
  "riskAssessment": {
    "staticKycScore": 45,
    "adjustedScore": 59,
    "sniHotFactor": 3,
    "riskLevel": "medium",
    "requiresEDD": false,
    "components": {
      "publicData": 12,
      "crossValidation": 18,
      "bookkeeping": 0
    },
    "triggers": {
      "isPEP": false,
      "hasSanctions": false,
      "hasBusinessProhibition": false,
      "vhMismatch": false,
      "addressMismatch": true,
      "missingFTax": false
    }
  },
  "nextStep": "/riskfragor/steg2"
}
\end{lstlisting}

\subsection{Beräkningsflöde}

\begin{enumerate}[noitemsep]
  \item Hämta SNI-kod från Bolagsverket $\rightarrow$ sniHotFactor (1–4)
  \item Beräkna \texttt{publicData} (0–30 poäng)
  \item Beräkna \texttt{crossValidation} (0–30 poäng)
  \item Sätt \texttt{bookkeeping} = 0 (fylls i senare)
  \item \texttt{staticKycScore} = publicData + crossValidation + bookkeeping
  \item \texttt{hotBonus} = (sniHotFactor - 1) $\times$ 10
  \item \texttt{adjustedScore} = staticKycScore + hotBonus
  \item Klassificera \texttt{riskLevel} enligt tröskeltabell (0-40 Low, 41-70 Medium, 71-90 High, >90 Critical)
  \item Kontrollera hårda triggers $\rightarrow$ eventuellt tvinga \texttt{critical}
  \item Kontrollera eskaleringstriggrar $\rightarrow$ eventuellt höj riskklass
\end{enumerate}

% ========================================================================
\section{Exempel: Restaurangföretag}
% ========================================================================

\subsection{Scenariodata}

\begin{itemize}[noitemsep]
  \item \textbf{Företag:} Celestial Restaurang AB
  \item \textbf{SNI-kod:} 56101 (Restaurangverksamhet) $\rightarrow$ Hot 3 (kontantintensiv)
  \item \textbf{Bolagsålder:} 2 år (registrerat 2023-01-15)
  \item \textbf{VD:} Bengt Bengtsson (personnr 197001011234)
  \item \textbf{Verkliga huvudmän:} Bengt 60\%, Lisa Larsson 40\%
  \item \textbf{PEP:} Nej
  \item \textbf{Sanktioner:} Nej
  \item \textbf{Avvikelser:} Adressmismatch (SPAR vs. Bolagsverket)
\end{itemize}

\subsection{Beräkning}

\textbf{Steg 1: Offentliga register}
\begin{itemize}[noitemsep]
  \item Bolagsålder 1–3 år: +4 poäng
  \item PEP: 0 poäng
  \item Sanktioner: 0 poäng
  \item Totalt: 4 poäng
\end{itemize}

\textbf{Steg 2: Korsvalidering}
\begin{itemize}[noitemsep]
  \item Adressmismatch: +8 poäng
  \item Övriga: 0 poäng
  \item Totalt: 8 poäng
\end{itemize}

\textbf{Steg 3: Bokföringsdata}
\begin{itemize}[noitemsep]
  \item Ännu ej analyserad: 0 poäng
\end{itemize}

\textbf{staticKycScore} = 4 + 8 + 0 = \textbf{12 poäng}

\textbf{Steg 4: Hot-bonus}
\begin{itemize}[noitemsep]
  \item SNI 56101 = Restaurang → Hot 3 (kontantintensiv)
  \item hotBonus = $(3-1) \times 10 = 20$ poäng
\end{itemize}

\textbf{adjustedScore} = $12 + 20 = \textbf{32 poäng}$

\textbf{Riskklass:} Low (0–40 poäng)

\textbf{requiresEDD:} false

\subsection{Kommentar}

Trots att restaurangbranschen är kontantintensiv (Hot 3, +20 poäng) är företaget fortfarande i \enquote{Low}-klass eftersom:
\begin{itemize}[noitemsep]
  \item Inga allvarliga avvikelser (endast mindre adressmismatch)
  \item Inga PEP/sanktioner/näringsförbud
  \item Bokföring ej analyserad ännu (kan höja score senare)
  \item Total score 32 poäng ligger under Low-gränsen (40 poäng)
\end{itemize}

Om senare SIE-analys visar hög kassaandel ($>40\%$) kan score stiga mot Medium-klass.

% ========================================================================
\section{Uppdatering vid bokföringsanalys}
% ========================================================================

När SIE-fil och bankkonto analyseras (senare i onboarding-flödet):

\begin{enumerate}[noitemsep]
  \item Beräkna \texttt{bookkeeping}-komponenten (0–40 poäng)
  \item Uppdatera \texttt{staticKycScore} = publicData + crossValidation + bookkeeping
  \item Räkna om \texttt{adjustedScore}
  \item Omklassificera \texttt{riskLevel}
  \item Om riskklass höjs: Notifiera PTL-ansvarig, kräv extra dokumentation
\end{enumerate}

\textbf{Exempel fortsättning (Restaurang):}

SIE-analys visar:
\begin{itemize}[noitemsep]
  \item Kassaandel 48\% av omsättning: +12 poäng
  \item Soliditet -5\%: 0 poäng (gräns är $< -20\%$)
  \item Totalt bookkeeping: 12 poäng
\end{itemize}

\textbf{Ny staticKycScore} = 4 + 8 + 12 = \textbf{24 poäng}

\textbf{Ny adjustedScore} = $24 + 20 = \textbf{44 poäng}$

\textbf{Ny riskklass:} Medium (41–70 poäng)

\textbf{Kommentar:} Med hög kassaandel eskaleras företaget från Low till Medium-klass. Extra dokumentation och fördjupad granskning av kontanthantering krävs nu enligt PTL.

% ========================================================================
\section{Sammanfattning och riktlinjer}
% ========================================================================

\subsection{Nyckelprinciper}

\begin{enumerate}[noitemsep]
  \item \textbf{Riskbaserat:} Högre risk = mer granskning, inte automatiskt avslag
  \item \textbf{Dynamiskt:} Score uppdateras när mer data tillkommer
  \item \textbf{Transparent:} Kunden ser sin riskklass och kan förklara avvikelser
  \item \textbf{Regelefterlevnad:} Baserat på PTL, FFFS 2009:1 och FATF-rekommendationer
\end{enumerate}

\subsection{Ansvar}

\begin{itemize}[noitemsep]
  \item \textbf{Systemet:} Automatisk beräkning av staticKycScore och riskLevel
  \item \textbf{PTL-ansvarig:} Godkänner High-klassade kunder, beslutar om EDD-åtgärder
  \item \textbf{VD:} Slutligt beslut vid Critical-klass (avslag eller anmälan)
\end{itemize}

\subsection{Dokumentation}

Alla riskbedömningar loggas i \texttt{event\_log.csv} (CSV-first approach) med:
\begin{itemize}[noitemsep]
  \item Timestamp
  \item onboardingId
  \item staticKycScore, adjustedScore, riskLevel
  \item Triggers som aktiverades
  \item PTL-ansvarigs kommentar (vid High/Critical)
\end{itemize}

\vspace{1cm}
\hrule
\vspace{0.5cm}
\textbf{Godkänd av:} [PTL-ansvarig namn] \\
\textbf{Datum:} \today \\
\textbf{Nästa revidering:} Årligen eller vid regeländringar

\end{document}
